% 猎户座（综述）
% license CCBYSA3
% type Wiki

本文根据 CC-BY-SA 协议转载翻译自维基百科\href{https://en.wikipedia.org/wiki/Orion_(constellation)}{相关文章}。

猎户座是在北天球冬季可见的一组显著恒星。它是现代划分的88个星座之一，同时也名列公元2世纪天文学家托勒密所列出的48个星座之中。其名称源自希腊神话中的一位猎人。

在北半球的冬季夜晚，猎户座最为醒目，另外还有五个星座与其一同构成“冬季六边形”星群。猎户座中最明亮的两颗恒星是参宿七（$\beta$）和参宿四（$\alpha$），二者均位列夜空中最亮的恒星之列；它们都是超巨星，并且亮度略有变化。此外，猎户座中还有六颗视星等亮于3.0的恒星，其中三颗构成了著名的猎户座腰带这一短而笔直的星群。猎户座还是一年一度猎户座流星雨的辐射点，该流星雨是与哈雷彗星相关的最强流星雨之一；同时，猎户座内还包含猎户座星云，这是夜空中最明亮的星云之一。
\subsection{特征}
\begin{figure}[ht]
\centering
\includegraphics[width=6cm]{./figures/c950af7d2bff8c26.png}
\caption{猎户座的星群示意图，如肉眼所见，并绘制了连线。} \label{fig_Orion_1}
\end{figure}
猎户座的西北方毗邻金牛座，西南方为波江座，南方为天兔座，东方为麒麟座，东北方为双子座。猎户座覆盖约594平方度，在88个星座中面积排名第26位。该星座的边界由比利时天文学家欧仁·德尔波特于1930年划定，其边界由一个26边形构成。在赤道坐标系中，这些边界的赤经范围介于 04\textsuperscript{h}43.3\textsuperscript{m} 与 06\textsuperscript{h}25.5\textsuperscript{m} 之间，而赤纬范围则介于 22.87° 与 −10.97° 之间。$^{[2]}$ 国际天文学联合会于1922年采用的该星座三字母缩写为 “Ori”。$^{[3]}$

猎户座在一月至四月的夜空中最为清晰可见，$^{[4]}$ 此时北半球正值冬季，而南半球为夏季。在热带地区（距赤道约8°以内），猎户座会通过天顶。

从五月到七月（北半球的夏季、南半球的冬季），猎户座位于白昼天空中，因此在大多数纬度地区不可见。然而，在南半球冬季的数月中，南极洲大部分地区即使在正午太阳也位于地平线以下。此时，恒星（以及猎户座，但仅限于最亮的恒星）会在地方正午前后数小时内于暮光中可见，出现在北方低空、天空中最明亮的区域，太阳正位于地平线下方不远处。在同一时刻、南极点本身（阿蒙森—斯科特南极站），参宿七仅高出地平线8°，而猎户座腰带几乎沿着地平线掠过。南半球的夏季月份中，猎户座在其他地区通常可于夜空中观测，但在南极洲却完全不可见，因为在南极圈以南的这一时期太阳并不落下。$^{[5][6]}$

在接近赤道的国家（例如肯尼亚、印度尼西亚、哥伦比亚、厄瓜多尔），猎户座在12月午夜前后位于头顶上方，而在2月则出现在傍晚的天空中。
\subsection{导航用途}
